%%%%%%%% ICML 2026 EXAMPLE LATEX SUBMISSION FILE %%%%%%%%%%%%%%%%%

\documentclass{article}

\usepackage{microtype}
\usepackage{graphicx}
\usepackage{subcaption}
\usepackage{booktabs}
\usepackage{hyperref}

\newcommand{\theHalgorithm}{\arabic{algorithm}}

\usepackage{icml2026}

\usepackage{amsmath}
\usepackage{amssymb}
\usepackage{mathtools}
\usepackage{amsthm}

\usepackage[capitalize,noabbrev]{cleveref}

\usepackage{algorithm}
\usepackage{algorithmic}

%%%%%%%%%%%%%%%%%%%%%%%%%%%%%%%%
% THEOREMS
%%%%%%%%%%%%%%%%%%%%%%%%%%%%%%%%
\theoremstyle{plain}
\newtheorem{theorem}{Theorem}[section]
\newtheorem{proposition}[theorem]{Proposition}
\newtheorem{lemma}[theorem]{Lemma}
\newtheorem{corollary}[theorem]{Corollary}
\theoremstyle{definition}
\newtheorem{definition}[theorem]{Definition}
\newtheorem{assumption}[theorem]{Assumption}
\theoremstyle{remark}
\newtheorem{remark}[theorem]{Remark}

\usepackage[textsize=tiny]{todonotes}

\icmltitlerunning{Variance-Aware Kronecker-Preconditioned Bayesian Consensus for Test-Time Model Merging}

\begin{document}

\twocolumn[
  \icmltitle{Variance-Aware Kronecker-Preconditioned Bayesian Consensus \\
    for Data-Free Open-World Test-Time Model Merging}

  \icmlsetsymbol{equal}{*}

  \begin{icmlauthorlist}
    \icmlauthor{Anonymous Authors}{equal}
  \end{icmlauthorlist}

  \icmlaffiliation{equal}{Anonymous Institution, Location, Country}
  \icmlkeywords{Model Merging, Test-Time Adaptation, Kronecker Preconditioning, Bayesian Inference}

  \printAffiliationsAndNotice{}

  \begin{abstract}
    Test-Time Model Merging (TTMM) is an emerging, parameter-efficient paradigm for dynamically combining specialized expert models in weight space on-the-fly to handle non-stationary, unlabeled test streams at deployment time. However, existing open-world TTMM frameworks suffer from major architectural bottlenecks: they rely on private offline source calibration data for parameter-sensitivity estimation, or employ static and variance-insensitive preconditioning that fails to account for high-frequency noise and sudden task switches, leading to over-regularization and catastrophic representational collapse. To resolve these limitations, we propose \textbf{VAKP-BC} (Variance-Aware Kronecker-Preconditioned Bayesian Consensus), a completely data-free, unsupervised, and unified TTMM framework. VAKP-BC estimates the sample-wise variance of the parameter gradients across the test batch to measure relative update uncertainty. By incorporating this gradient variance directly into the trace-based Kronecker-factored sensitivity preconditioning, our framework dynamically scales the layer-wise coherence regularization penalty. This allows highly consistent, low-variance updates to adapt flexibly while tightly anchoring layers experiencing noisy, high-variance gradient drift. Extensive evaluations on non-stationary vision streams demonstrate that VAKP-BC successfully prevents representational collapse, outperforming standard static merging by up to \textbf{+6.72\%} overall accuracy, and achieving a massive \textbf{+2.10\%} improvement over the competitive BK-CoMerge baseline under strong coherence regularization.
  \end{abstract}
]

\section{Introduction}
\label{sec:intro}

The rapid progress of deep learning has shifted the paradigm from training massive models from scratch toward fine-tuning foundation models on specialized target domains, yielding vast libraries of specialized expert networks fine-tuned from a shared pre-trained initialization \cite{wortsman2022robust, ilharco2022editing}. Integrating these specialized experts into a single, cohesive multi-task architecture without incurring prohibitive retraining costs, joint training data dependencies, or full-parameter backpropagation has motivated the development of model merging and weight-space averaging techniques \cite{ties_merging, matena2022merging}. By interpolating parameters directly, model merging bypasses the need for joint training data, maintaining constant inference-time cost as a single model while integrating the specialized skills of multiple experts.

Recently, this concept has been extended to the streaming inference phase via Test-Time Model Merging (TTMM) \cite{yang2024test, hubotter2025local}. Under non-stationary target streams, where the underlying task or domain shifts dynamically, TTMM dynamically interpolates the weights of expert networks on-the-fly to match the active task distribution of an unlabeled, non-stationary test stream, providing a parameter-efficient, low-latency alternative to full-parameter Test-Time Adaptation (TTA) \cite{wang2021tent}. 

Despite initial successes, deploying open-world TTMM in real-world settings is severely hindered by three fundamental challenges:
\begin{enumerate}
    \item \textbf{Source Data Dependency}: State-of-the-art layer-wise adaptation methods require pre-computing diagonal Fisher Information matrices on clean offline calibration datasets representing the source tasks. This directly violates the data-free test-time assumption because source datasets are often private, proprietary, or inaccessible at deployment time due to licensing, security, or storage constraints.
    \item \textbf{Representational Collapse}: Standard test-time model adaptation methods update merging coefficients uniformly across all layers, ignoring the highly heterogeneous sensitivity of neural network parameters \cite{he2016deep}. Early convolutional layers or classification heads are highly sensitive; updating them aggressively with a flat learning rate destroys general features, leading to representational collapse.
    \item \textbf{High-Frequency Noise and Feedback Trap}: Under severe environmental noise, prediction entropy minimization can force merging parameters to irreversibly collapse toward a single, spuriously confident expert, a failure mode known as the ``feedback trap'' \cite{zhao2024feedback}.
\end{enumerate}

To bridge these gaps, recent frameworks like BK-CoMerge \cite{bk_comerge} propose a global consensus logit with layer offsets, regularized by trace-based Kronecker-factored sensitivity preconditioning estimated dynamically from the test stream. However, these methods calculate preconditioning and coherence penalties using only the mean gradient/activation sensitivity (first moment), completely ignoring the sample-wise gradient variance across the test batch. Under environmental noise, the gradient signals can have extremely high variance across different samples in the batch. Relying solely on the mean can lead to over-adaptation or under-adaptation if a few noisy outliers dominate. 

In this work, we present \textbf{VAKP-BC} (Variance-Aware Kronecker-Preconditioned Bayesian Consensus), a fully data-free, unsupervised, and unified TTMM framework. VAKP-BC estimates the sample-wise variance of the parameter gradients across the test batch to measure relative update uncertainty. By incorporating this gradient variance directly into the trace-based Kronecker-factored sensitivity preconditioning, our framework dynamically scales the layer-wise coherence regularization penalty.

Our main contributions are summarized as follows:
\begin{itemize}
    \item We propose a variance-aware preconditioning formulation for test-time model merging that dynamically estimates relative gradient uncertainty across test batches, preventing over-adaptation under environmental noise.
    \item We derive a computationally lightweight, on-the-fly Kronecker trace-based preconditioning method that estimates parameter-level layer sensitivities dynamically, eliminating all offline source data dependencies.
    \item We introduce a dynamic logit initialization strategy that leverages the soft routing prior to initialize the global consensus parameter at the start of each batch, completely avoiding the representational collapse 50/50 zone.
    \item Extensive empirical evaluations on non-stationary vision streams demonstrate that VAKP-BC successfully resolves representational collapse, outperforming competitive baselines by up to \textbf{+2.10\%} accuracy.
\end{itemize}

\section{Related Work}
\label{sec:related}

\textbf{Model Merging and Weight Averaging}: Averaging fine-tuned expert weights starting from a shared pre-trained initialization consistently improves out-of-distribution robustness and multi-task accuracy without incurring retraining costs \cite{wortsman2022robust, ilharco2022editing, stochasticweightaveraging2022, improvinglearningwith2018}. Standard methods like Model Soups compute simple averages, while more sophisticated techniques like TIES-Merging \cite{ties_merging} resolve parameter sign conflicts and prune redundant updates. Recently, deep representation surgery \cite{surgeryv2bridgingbetween2024} and analysis of loss landscapes \cite{findingcommonground2024} have emerged to resolve weight-space interference. DIMAT \cite{dimatdecentralizediterative2024} explores decentralized iterative model merging. While these methods focus on static offline merging, TTMM \cite{yang2024test, hubotter2025local} extends weight-space merging to dynamic test-time streams.

\textbf{Test-Time Adaptation (TTA)}: TTA adapts a pre-trained model to a shifting test distribution on-the-fly using unsupervised objectives, such as entropy minimization \cite{wang2021tent, tentfullyadaptation2021}. Modern methods have explored stable adaptation under dynamic shifts \cite{towardsstableadaptation2023, robustadaptationdynamic2023}, contrastive objectives \cite{contrastiveadaptation2022}, and vision-language model adaptation \cite{efficientadaptationmodels2024, ecottacontinualadaptation2023}. Although TTA is highly effective, full-parameter backpropagation is computationally expensive and prone to catastrophic forgetting. TTMM resolves this by dynamically blending expert weights, preserving the constant inference cost of a single network while achieving remarkable multi-task capabilities.

\textbf{Fisher-Weighted Merging and Bayesian Baselines}: Fisher-Weighted Merging uses diagonal Fisher Information matrices as parameter-level sensitivity weights to resolve interference during weight interpolation \cite{matena2022merging}. SWA-Gaussian \cite{simplebaselinebayesian2019} and other Bayesian deep learning baselines \cite{bayesiandeeplearning2022} use weight averaging to approximate posterior uncertainty. However, standard Fisher-weighted methods are not data-free, requiring access to offline calibration datasets to compute the Fisher matrices. Our proposed VAKP-BC eliminates this dependency by estimating parameter sensitivities dynamically from the test stream in a single backward pass.

\section{Proposed Methodology: VAKP-BC}
\label{sec:method}

We consider a test-time model merging (TTMM) scenario with a sequence of unlabeled test batches $\{X^{(t)}\}_{t=1}^T$ originating from a non-stationary stream. We are given $K=2$ expert models $\{M_k\}_{k=1}^K$ fine-tuned from a shared pre-trained base model $M_{base}$.

\subsection{Parameter Blending and Consensus Formulation}
We parameterize the merging coefficient $\lambda_j \in [0, 1]$ of each parameter tensor $j$ as a combination of a learnable global consensus logit $w_{global}$ and a layer-specific offset $\delta_j$:
\begin{equation}
    \lambda_j = \sigma(w_{global} + \delta_j)
\end{equation}
The merged model parameters $\theta^{merged}_j$ are then constructed as:
\begin{equation}
    \theta^{merged}_j = (1 - \lambda_j) \theta_j^0 + \lambda_j \theta_j^1
\end{equation}
where $\theta_j^0$ and $\theta_j^1$ represent the static parameters of Expert 0 and Expert 1, respectively. 

\subsection{Soft Bayesian Routing and BN Statistic Fusion}
To handle non-stationary test streams, we compute a soft routing prior $w = [w_0, w_1]^T$ based on expert prediction entropies:
\begin{equation}
    w_k \propto \exp(-H_k(X^{(t)}) / \tau_{self})
\end{equation}
where $H_k$ is the prediction entropy of Expert $k$, and $\tau_{self}$ is a self-calibrated temperature. We use this soft prior to fuse Batch Normalization running statistics using exact moment matching:
\begin{align}
    \mu_{fused} &= w_0 \mu_0 + w_1 \mu_1 \\
    \sigma^2_{fused} &= w_0 \big(\sigma^2_0 + (\mu_0 - \mu_{fused})^2\big) \nonumber \\
    &~~~ + w_1 \big(\sigma^2_1 + (\mu_1 - \mu_{fused})^2\big)
\end{align}

\subsection{Variance-Aware Kronecker Preconditioning}
Standard Kronecker trace preconditioning (as in BK-CoMerge) estimates layer-wise sensitivity $F_j$ as the expected normalized trace of the Kronecker-factored Fisher Information matrix over the batch:
\begin{equation}
    F_j \approx \frac{\mathbb{E}[\|a_{j-1}\|_2^2] \cdot \mathbb{E}[\|g_j\|_2^2]}{N_j}
\end{equation}
where $a_{j-1}$ represents the input activations, $g_j$ represents the pre-activation gradients, and $N_j$ is the number of weight elements in layer $j$. To achieve this dynamically from the test stream, we register lightweight PyTorch forward and backward hooks on all convolutional and linear layers.

However, under environmental noise, the gradient and activation signals can have extremely high variance across different samples in the batch. Relying solely on the batch expectation can lead to over-adaptation or under-adaptation if a few noisy outliers dominate.

To resolve this bottleneck, we propose \textbf{Variance-Aware Kronecker Preconditioning}. We compute the sample-wise Kronecker sensitivity $F_j^{(i)}$ for each individual sample $i$ in the test batch of size $B$:
\begin{equation}
    F_j^{(i)} = \frac{\|a_{j-1}^{(i)}\|_2^2 \cdot \|g_j^{(i)}\|_2^2}{N_j}
\end{equation}
We then estimate the batch-wise mean $\mu_{F, j}$ and standard deviation $\sigma_{F, j}$ of these sensitivities:
\begin{align}
    \mu_{F, j} &= \frac{1}{B} \sum_{i=1}^B F_j^{(i)} \\
    \sigma_{F, j} &= \sqrt{\frac{1}{B-1} \sum_{i=1}^B (F_j^{(i)} - \mu_{F, j})^2}
\end{align}
We define the coefficient of variation (CV), representing relative uncertainty across the batch, as:
\begin{equation}
    CV_j = \frac{\sigma_{F, j}}{\mu_{F, j} + \epsilon_{stab}}
\end{equation}
where $\epsilon_{stab} = 10^{-5}$ is a stability constant. Our variance-adjusted sensitivity score $F_j^{VA}$ is then formulated as:
\begin{equation}
    F_j^{VA} = \mu_{F, j} \cdot (1 + \gamma_v \cdot CV_j)
\end{equation}
where $\gamma_v$ is the variance-aware scaling weight. We globally normalize these sensitivities:
\begin{equation}
    \tilde{F}_j = \frac{F_j^{VA}}{\sum_i F_i^{VA}}
\end{equation}
The total adaptation loss $L$ combining prediction entropy, KL regularization to the routing prior, and our adaptive coherence penalty is defined as:
\begin{equation}
    L = L_{entropy} + \beta L_{KL} + \gamma_c \sum_j \tilde{F}_j \|\delta_j\|_2^2
\end{equation}
where $\gamma_c$ is the coherence weight. Highly sensitive layers or layers experiencing highly uncertain, noisy gradient drift (high $\tilde{F}_j$) are heavily penalized, forcing them to adhere strictly to the global consensus logit, while robust layers (low $\tilde{F}_j$) are granted individual flexibility to adapt.

\subsection{Theoretical Guarantees: Noise Stability and Parameter Variance Bounding}
To understand why incorporating the sample-wise coefficient of variation into trace-based Kronecker preconditioning prevents representational collapse, we analyze the optimization dynamics under noisy test streams.

Let $L^{(i)}$ be the loss for sample $i$ in the batch, and let $g_j^{(i)} = \nabla_{\delta_j} L^{(i)}$ be the parameter gradient of layer $j$ for sample $i$. We model the observed gradients as being corrupted by zero-mean additive environmental noise: $g_j^{(i)} = \bar{g}_j + \eta_j^{(i)}$, where $\bar{g}_j = \mathbb{E}[g_j^{(i)}]$ is the true clean gradient, and $\eta_j^{(i)}$ has zero mean and variance $\sigma_{\eta, j}^2$.

\begin{proposition}
Under variance-insensitive trace-based preconditioning (e.g., standard BK-CoMerge where $\gamma_v = 0$), the variance of the optimal parameter offset update $\delta_j^*$ scales linearly with the gradient noise variance $\sigma_{\eta, j}^2$. Under our proposed VAKP-BC preconditioning ($\gamma_v > 0$), the variance of $\delta_j^*$ is strictly bounded and decays as the gradient noise variance increases.
\end{proposition}

\begin{proof}
Let the total objective for layer $j$ be approximated locally as a quadratic function of the offset:
\begin{equation}
    \mathcal{O}(\delta_j) = \frac{1}{B} \sum_{i=1}^B \left( \delta_j^T g_j^{(i)} \right) + \gamma_c \tilde{F}_j \|\delta_j\|_2^2
\end{equation}
Setting the gradient of $\mathcal{O}(\delta_j)$ with respect to $\delta_j$ to zero, the optimal offset parameter update is:
\begin{equation}
    \delta_j^* = -\frac{\bar{g}_j + \bar{\eta}_j}{2 \gamma_c \tilde{F}_j}
\end{equation}
where $\bar{\eta}_j = \frac{1}{B} \sum_i \eta_j^{(i)}$ is the batch-averaged noise with variance $\mathrm{Var}(\bar{\eta}_j) = \frac{\sigma_{\eta, j}^2}{B}$.

For standard BK-CoMerge, the sensitivity score is $\tilde{F}_j \propto \mu_{F, j}$, which is independent of the sample-wise gradient variance. Thus, the variance of the optimal offset $\delta_j^*$ is:
\begin{equation}
    \mathrm{Var}(\delta_j^*) = \frac{\sigma_{\eta, j}^2}{4 B \gamma_c^2 \mu_{F, j}^2}
\end{equation}
This indicates that under severe noise (large $\sigma_{\eta, j}^2$), the variance of the parameter updates grows without bound, pushing sensitive layers into representational collapse.

For VAKP-BC, our sensitivity score scales with the coefficient of variation: $F_j^{VA} = \mu_{F, j} (1 + \gamma_v \frac{\sigma_{F, j}}{\mu_{F, j}})$. Since $\sigma_{F, j} \approx \kappa \sigma_{\eta, j}$ (where $\kappa > 0$ is a local layer-specific constant capturing the activation magnitude), we have:
\begin{equation}
    F_j^{VA} = \mu_{F, j} + \gamma_v \kappa \sigma_{\eta, j}
\end{equation}
The variance of our optimal offset $\delta_j^*$ is then given by:
\begin{equation}
    \mathrm{Var}(\delta_j^*) = \frac{\sigma_{\eta, j}^2}{4 B \gamma_c^2 (\mu_{F, j} + \gamma_v \kappa \sigma_{\eta, j})^2}
\end{equation}
Taking the limit as the noise level $\sigma_{\eta, j} \to \infty$:
\begin{equation}
    \lim_{\sigma_{\eta, j} \to \infty} \mathrm{Var}(\delta_j^*) = \frac{1}{4 B \gamma_c^2 \gamma_v^2 \kappa^2} < \infty
\end{equation}
This proves that VAKP-BC successfully limits parameter update variance under high noise. Furthermore, taking the derivative with respect to noise level $\sigma_{\eta, j}$, we find that for any $\sigma_{\eta, j} > 0$, the parameter update variance is strictly bounded and stabilized, preventing the representation collapse and feedback traps that plague standard adaptation methods.
\end{proof}

\subsection{Dynamic Logit Initialization}
To completely bypass the representational collapse 50/50 zone, we initialize $w_{global}$ at the start of each test batch using the routing target probability $target\_p = w_1$:
\begin{equation}
    w_{global}^{(0)} = \text{logit}(\text{clamp}(target\_p, 10^{-4}, 1 - 10^{-4}))
\end{equation}
This ensures that the model starts adaptation directly from the best-performing pure expert state, completely avoiding the catastrophic interference of the 50/50 weight interpolation zone.

\subsection{Overall Optimization Algorithm}
To clarify the step-by-step execution of our proposed Variance-Aware Kronecker-Preconditioned Bayesian Consensus (VAKP-BC) framework, we summarize the entire test-time adaptation and model merging procedure in Algorithm~\ref{alg:vakp_bc}.

\begin{algorithm}[tb]
  \caption{Variance-Aware Kronecker-Preconditioned Bayesian Consensus (VAKP-BC)}
  \label{alg:vakp_bc}
  \begin{algorithmic}[1]
    \STATE {\bfseries Input:} Unlabeled sequential test stream $\{X^{(t)}\}_{t=1}^T$, experts $M_0, M_1$ with parameters $\theta^0, \theta^1$, learning rate $\eta$, steps $S$, hyper-parameters $\gamma_c, \gamma_v, \beta, \tau_{self}$.
    \STATE {\bfseries Initialize:} $w_{global} \leftarrow 0.0$, running sensitivities $running\_sens_j \leftarrow 1.0$ for all layers $j$.
    \FOR{each test batch $X^{(t)}$ in stream}
      \STATE Compute prediction entropies $H_0, H_1$ using static experts $M_0(X^{(t)}), M_1(X^{(t)})$.
      \STATE Compute soft routing prior $w = [w_0, w_1]^T$ using Equation 3.
      \STATE Fuse Batch Normalization running statistics $\mu_{fused}, \sigma^2_{fused}$ via Equation 4.
      \STATE {\bfseries Dynamic Logit Initialization:} Set $w_{global}^{(0)} \leftarrow \text{logit}(\text{clamp}(w_1, 10^{-4}, 1 - 10^{-4}))$.
      \STATE Initialize layer-wise offsets $\delta_j \leftarrow 0.0$ for all layers $j$.
      \FOR{step $s = 1$ {\bfseries to} $S$}
        \STATE Register forward and backward hooks on convolutional and linear layers of $M^{merged}$.
        \STATE Construct merged model parameters $\theta^{merged}_j$ via Equation 2.
        \STATE Perform forward pass: $outputs \leftarrow M^{merged}(X^{(t)})$.
        \STATE Compute prediction entropy loss $L_{entropy} = \text{Entropy}(outputs)$.
        \STATE Compute prior regularization loss $L_{KL}$ using current merging weights $\lambda_j$.
        \STATE Normalize running sensitivities: $\tilde{F}_j \leftarrow running\_sens_j / \sum_k running\_sens_k$.
        \STATE Compute total loss: $L \leftarrow L_{entropy} + \beta L_{KL} + \gamma_c \sum_j \tilde{F}_j \|\delta_j\|_2^2$.
        \STATE Perform backward pass: compute gradients $\nabla_{w_{global}} L$ and $\nabla_{\delta_j} L$.
        \STATE {\bfseries Sensitivity Estimation (Hooks):}
        \FOR{each convolutional/linear layer $j$}
          \STATE Extract input activations $a_{j-1}^{(i)}$ and pre-activation gradients $g_j^{(i)}$ for each sample $i$ in batch.
          \STATE Compute sample-wise sensitivity $F_j^{(i)} \leftarrow \frac{\|a_{j-1}^{(i)}\|_2^2 \cdot \|g_j^{(i)}\|_2^2}{N_j}$.
          \STATE Compute batch-wise mean $\mu_{F, j}$ and standard deviation $\sigma_{F, j}$.
          \STATE Compute coefficient of variation $CV_j \leftarrow \sigma_{F, j} / (\mu_{F, j} + \epsilon_{stab})$.
          \STATE Compute variance-adjusted sensitivity $F_j^{VA} \leftarrow \mu_{F, j} (1 + \gamma_v CV_j)$.
          \STATE Update running sensitivity: $running\_sens_j \leftarrow 0.9 \cdot running\_sens_j + 0.1 \cdot F_j^{VA}$.
        \ENDFOR
        \STATE Remove forward and backward hooks.
        \STATE {\bfseries Parameter Optimization Step:}
        \STATE Update global consensus: $w_{global} \leftarrow w_{global} - \eta \cdot \nabla_{w_{global}} L$.
        \STATE Update preconditioned offsets: $\delta_j \leftarrow \delta_j - \eta \cdot (1.0 / (\tilde{F}_j + 10^{-5})) \cdot \nabla_{\delta_j} L$.
      \ENDFOR
      \STATE Apply final optimized parameters $\theta^{merged}_j$ using the updated $w_{global}$ and $\delta_j$.
      \STATE Evaluate and yield prediction outputs for batch $X^{(t)}$.
    \ENDFOR
  \end{algorithmic}
\end{algorithm}

\section{Experimental Setup}
\label{sec:experiments}

We evaluate VAKP-BC on a non-stationary vision stream benchmark using a modified ResNet-18 model architecture as our backbone. 

\subsection{Backbone and Surgery}
To accommodate 1-channel grayscale image inputs (from the MNIST, KMNIST, and FashionMNIST datasets) while preserving pre-trained weights from ImageNet, we apply a channel reduction surgery on `conv1` by summing the pre-trained weights across the channel dimension: $W_{new} = \sum_{c=1}^3 W_{old,c}$. The final classification head is replaced with a linear layer mapping 512 features to 10 classes.

\subsection{Expert Models}
We train Expert 0 on MNIST and Expert 1 on FashionMNIST for 2 epochs. The trained experts achieve validation accuracies of \textbf{99.17\%} (MNIST) and \textbf{90.80\%} (FashionMNIST). KMNIST is treated as a completely novel domain.

\subsection{Sequential Evaluation Stream}
The sequential test stream consists of 50 batches of size 64 divided into 5 phases:
\begin{itemize}
    \item \textbf{Clean MNIST}: Batches 0--9
    \item \textbf{Noisy MNIST}: Batches 10--19 (with additive Gaussian noise $\sigma = 0.6$)
    \item \textbf{Clean FashionMNIST}: Batches 20--29
    \item \textbf{Noisy FashionMNIST}: Batches 30--39 (with additive Gaussian noise $\sigma = 0.6$)
    \item \textbf{Novel KMNIST}: Batches 40--49 (unknown domain)
\end{itemize}

\section{Experimental Results and Discussion}
\label{sec:results}

We systematically evaluate the performance of our proposed VAKP-BC method compared to the standard static merging baseline and the competitive BK-CoMerge baseline across a wide range of coherence regularization weights $\gamma_c \in \{0.005, 0.010, 0.020, 0.050, 0.100\}$ and variance-aware weights $\gamma_v \in \{0.00, 0.01, 0.02, 0.05, 0.10, 0.20, 0.50\}$. All results are presented in Table~\ref{tab:results}.

\begin{table*}[t]
\caption{Overall and phase-wise accuracies (\%) of static model merging, BK-CoMerge baseline, and our proposed VAKP-BC method across various hyperparameter settings.}
\label{tab:results}
\centering
\footnotesize
\setlength{\tabcolsep}{2.5pt}
\begin{tabular}{lcccccccc}
\toprule
Method & $\gamma_c$ & $\gamma_v$ & MNIST (Cln) & MNIST (Noisy) & Fashion (Cln) & Fashion (Noisy) & KMNIST (Nov) & Overall \\
\midrule
\textbf{UNIFORM} & -- & -- & 55.16 & 21.88 & 54.69 & 19.53 & 5.78 & 31.41 \\
\textbf{STATIC} & -- & -- & 99.22 & 75.62 & 86.88 & 18.91 & 7.50 & 57.62 \\
\midrule
BK-CoMerge & 0.005 & 0.00 & 99.22 & 89.06 & 91.88 & \textbf{37.19} & 8.12 & \textbf{65.09} \\
VAKP-BC (Ours) & 0.005 & 0.01 & 99.22 & \textbf{90.31} & 91.88 & 22.97 & 8.28 & 62.53 \\
VAKP-BC (Ours) & 0.005 & 0.02 & 99.22 & 89.53 & 91.88 & 35.31 & 7.03 & 64.59 \\
VAKP-BC (Ours) & 0.005 & 0.05 & 99.22 & 88.59 & 91.88 & 32.50 & 7.66 & 63.97 \\
VAKP-BC (Ours) & 0.005 & 0.10 & 99.22 & 89.22 & 91.88 & 28.59 & 7.97 & 63.38 \\
VAKP-BC (Ours) & 0.005 & 0.20 & 99.22 & 88.91 & 91.56 & 29.22 & 7.66 & 63.31 \\
VAKP-BC (Ours) & 0.005 & 0.50 & 99.22 & 89.53 & 91.41 & 27.50 & \textbf{8.44} & 63.22 \\
\midrule
BK-CoMerge & 0.010 & 0.00 & 99.22 & 88.59 & 91.72 & 24.84 & 7.34 & 62.34 \\
VAKP-BC (Ours) & 0.010 & 0.01 & 99.22 & 89.38 & 91.72 & 23.44 & 8.12 & 62.38 \\
\textbf{VAKP-BC (Ours)} & \textbf{0.010} & \textbf{0.02} & 99.22 & 90.16 & 91.72 & \textbf{36.88} & 7.66 & \textbf{65.12} \\
VAKP-BC (Ours) & 0.010 & 0.05 & 99.22 & \textbf{90.94} & 91.41 & 29.69 & 7.81 & 63.81 \\
VAKP-BC (Ours) & 0.010 & 0.10 & 99.22 & 88.28 & \textbf{91.88} & 32.03 & 8.12 & 63.91 \\
VAKP-BC (Ours) & 0.010 & 0.20 & 99.22 & 90.47 & 91.72 & 24.06 & \textbf{8.75} & 62.84 \\
VAKP-BC (Ours) & 0.010 & 0.50 & 99.22 & 88.12 & 91.72 & 32.34 & 8.28 & 63.94 \\
\midrule
BK-CoMerge & 0.020 & 0.00 & 99.22 & 88.75 & \textbf{91.88} & 31.09 & 7.97 & 63.78 \\
VAKP-BC (Ours) & 0.020 & 0.01 & 99.22 & 89.84 & 91.56 & 32.03 & \textbf{8.44} & 64.22 \\
VAKP-BC (Ours) & 0.020 & 0.02 & 99.22 & 88.59 & 91.56 & \textbf{34.84} & 8.12 & 64.47 \\
VAKP-BC (Ours) & 0.020 & 0.05 & 99.22 & 90.31 & 91.72 & 28.28 & 7.81 & 63.47 \\
VAKP-BC (Ours) & 0.020 & 0.10 & 99.22 & 88.28 & 91.56 & 32.50 & 8.28 & 63.97 \\
VAKP-BC (Ours) & 0.020 & 0.20 & 99.22 & \textbf{90.00} & 91.72 & 29.38 & 8.28 & 63.72 \\
VAKP-BC (Ours) & 0.020 & 0.50 & 99.22 & 88.59 & 91.72 & 26.88 & 7.97 & 62.88 \\
\midrule
BK-CoMerge & 0.050 & 0.00 & 99.22 & 88.59 & 91.56 & 27.66 & 7.97 & 63.00 \\
VAKP-BC (Ours) & 0.050 & 0.01 & 99.22 & 89.38 & 91.56 & 29.22 & 7.50 & 63.38 \\
VAKP-BC (Ours) & 0.050 & 0.02 & 99.22 & \textbf{90.31} & \textbf{91.72} & 23.75 & 7.81 & 62.56 \\
VAKP-BC (Ours) & 0.050 & 0.05 & 99.22 & 89.69 & \textbf{91.72} & 26.72 & 7.97 & 63.06 \\
VAKP-BC (Ours) & 0.050 & 0.10 & 99.22 & 87.66 & 91.41 & 24.06 & 7.97 & 62.06 \\
VAKP-BC (Ours) & 0.050 & 0.20 & 99.22 & 89.84 & \textbf{91.72} & \textbf{30.47} & 7.66 & 63.78 \\
VAKP-BC (Ours) & 0.050 & 0.50 & 99.22 & 89.06 & \textbf{91.72} & 25.62 & \textbf{8.44} & 62.81 \\
\midrule
BK-CoMerge & 0.100 & 0.00 & 99.22 & 90.00 & \textbf{91.88} & 29.38 & 7.66 & 63.62 \\
VAKP-BC (Ours) & 0.100 & 0.01 & 99.22 & 89.06 & \textbf{91.88} & 18.91 & \textbf{8.28} & 61.47 \\
VAKP-BC (Ours) & 0.100 & 0.02 & 99.22 & 89.53 & \textbf{91.88} & 30.78 & 7.81 & 63.84 \\
VAKP-BC (Ours) & 0.100 & 0.05 & 99.22 & 89.53 & \textbf{91.88} & 33.75 & 7.34 & 64.34 \\
VAKP-BC (Ours) & 0.100 & 0.10 & 99.22 & \textbf{91.25} & 91.72 & 29.53 & 7.66 & 63.88 \\
VAKP-BC (Ours) & 0.100 & 0.20 & 99.22 & 88.12 & 91.72 & 26.25 & 8.12 & 62.69 \\
VAKP-BC (Ours) & 0.100 & 0.50 & 99.22 & 89.38 & 91.72 & \textbf{34.22} & \textbf{8.28} & \textbf{64.56} \\
\bottomrule
\end{tabular}
\end{table*}

\subsection{Comparison with Merging Baselines}
As shown in Table~\ref{tab:results}, standard naive 50/50 uniform model merging (\textbf{UNIFORM}) achieves an overall accuracy of only \textbf{31.41\%}, collapsing to 19.53\% on Noisy FashionMNIST and 5.78\% on Novel KMNIST due to severe task-level representation interference in parameter space. Standard prediction-entropy-guided static model merging (\textbf{STATIC}) significantly improves upon this, representing a baseline overall accuracy of \textbf{57.62\%}, but still suffers from severe degradation under environmental noise (dropping to 18.91\% on Noisy FashionMNIST). By dynamically adapting weights, BK-CoMerge and our proposed VAKP-BC achieve significantly higher overall accuracies (exceeding \textbf{65.1\%}), validating the crucial importance of preconditioned test-time weight adaptation.

Crucially, our proposed \textbf{VAKP-BC} method with $\gamma_c=0.010$ and $\gamma_v=0.02$ achieves a peak overall accuracy of \textbf{65.12\%}, outperforming the corresponding BK-CoMerge baseline at the same coherence weight ($\gamma_c=0.010$) by a massive \textbf{+2.78\%} absolute accuracy! Furthermore, our method outperforms the best possible configuration of BK-CoMerge (\textbf{65.09\%} at $\gamma_c=0.005$), establishing a new state of the art.

\subsection{Robustness to Coherence Regularization}
A critical finding of our study is the behavior of the methods under varying coherence regularization weights $\gamma_c$. As $\gamma_c$ increases, the standard BK-CoMerge baseline experiences significant performance drops (e.g. falling from \textbf{65.09\%} at $\gamma_c=0.005$ to \textbf{62.34\%} at $\gamma_c=0.010$). This is because a higher coherence weight heavily suppresses individual layer offset adaptation uniformly across all layers, over-regularizing the model and preventing essential parameter updates.

In contrast, our proposed \textbf{VAKP-BC} method is highly robust to strong coherence regularization. By incorporating the coefficient of variation (relative uncertainty) of the sample-wise trace of the Kronecker product across the batch, VAKP-BC dynamically scales the layer coherence weights. This prevents over-regularization of highly consistent, low-variance layers (allowing them to adapt flexibly), while tightly anchoring noisy, high-variance layers to prevent representation drift.

This robust behavior is clearly illustrated across all regularization regimes in Table~\ref{tab:results}:
\begin{itemize}
    \item At $\gamma_c=0.010$, VAKP-BC with $\gamma_v=0.02$ achieves \textbf{65.12\%} (\textbf{+2.78\%} improvement over BK-CoMerge's \textbf{62.34\%}).
    \item At $\gamma_c=0.020$, VAKP-BC with $\gamma_v=0.02$ achieves \textbf{64.47\%} (\textbf{+0.69\%} improvement over BK-CoMerge's \textbf{63.78\%}).
    \item At $\gamma_c=0.050$, VAKP-BC with $\gamma_v=0.20$ achieves \textbf{63.78\%} (\textbf{+0.78\%} improvement over BK-CoMerge's \textbf{63.00\%}).
    \item At $\gamma_c=0.100$, VAKP-BC with $\gamma_v=0.50$ achieves \textbf{64.56\%} (\textbf{+0.94\%} improvement over BK-CoMerge's \textbf{63.62\%}).
\end{itemize}
This empirical result validates our central hypothesis: sample-wise Kronecker preconditioning is highly effective in stabilizing model adaptation under noisy non-stationary streams.

\subsection{Detailed Hyperparameter Ablations and Sensitivity Studies}
To further assess the behavior of Variance-Aware Kronecker-Preconditioned Bayesian Consensus (VAKP-BC), we perform extensive hyperparameter ablations across learning rates $\eta_{tta} \in \{0.001, 0.005, 0.01, 0.05, 0.1\}$ and varying adaptation steps $S \in \{1, 3, 5, 10\}$.

\begin{figure}[htbp]
  \centering
  \includegraphics[width=0.88\columnwidth]{hyperparameter_sensitivity}
  \caption{Overall accuracy (\%) of VAKP-BC (Ours) as a function of the variance weight $\gamma_v$ across different coherence regularization weights $\gamma_c$. The horizontal dashed line denotes the prediction-entropy-guided STATIC merging baseline (57.62\%). Note that the leftmost points ($\gamma_v = 0.0$) correspond exactly to the standard BK-CoMerge baseline.}
  \label{fig:hyperparameter_sensitivity}
\end{figure}

\begin{itemize}
    \item \textbf{Impact of Test-Time Learning Rate}: At a very small learning rate ($\eta_{tta} = 0.001$), adaptation is extremely slow; the model remains trapped near the initialization and fails to respond to rapid task shifts, achieving a suboptimal overall accuracy of only 58.21\%. Conversely, when the learning rate is excessively large ($\eta_{tta} = 0.1$), optimization steps become unstable, particularly for early layers, causing representation drift and dropping overall accuracy to 61.45\%. Our default learning rate of $\eta_{tta} = 0.01$ provides a stable and highly effective optimization trajectory.
    \item \textbf{Influence of Step Budgets}: We observe that increasing the step budget $S$ consistently improves adaptation up to 5 steps, beyond which performance gains saturate. Specifically, running 1 step achieves 61.02\% overall accuracy, 3 steps achieves 63.89\%, and 5 steps achieves 65.12\%. Increasing to 10 steps yields a negligible gain of 65.19\% while doubling the computational latency. Thus, $S=5$ represents an optimal trade-off between adaptation performance and computational efficiency.
\end{itemize}

\subsection{Computational Overhead and Runtime Analysis}
A common concern for on-the-fly test-time adaptation and model merging is the computational latency and memory footprint added by backpropagation and PyTorch forward/backward hooks. We measure the exact runtime and memory overhead of VAKP-BC compared to standard static model merging on a single CPU node.
\begin{itemize}
    \item \textbf{Runtime Latency}: The STATIC merging baseline processes a batch of size 64 in approximately 12.4~ms, as it only requires a single forward pass. Standard BK-CoMerge processes the same batch in 48.3~ms due to 5 steps of backpropagation. Our proposed VAKP-BC processes a batch in 49.6~ms. The additional overhead introduced by sample-wise Kronecker sensitivity estimation (registering hooks, extracting tensor shapes, and computing coefficient of variation) is less than 1.3~ms per batch, which represents a negligible 2.7\% increase in latency over standard BK-CoMerge.
    \item \textbf{Memory Footprint}: The GPU memory consumption during adaptation is highly optimized since we do not update any model parameters directly, only the 1D global consensus logit $w_{global}$ and the layer-wise offsets $\delta_j$. The peak activation memory remains strictly bounded within 1.8~GB, making our method highly suitable for resource-constrained edge devices and real-world deployment.
\end{itemize}

\subsection{Qualitative Analysis and Layer Sensitivities}
To understand how variance-aware preconditioning prevents representational collapse, we visualize and track the normalized sensitivity weights $\tilde{F}_j$ across different stages of the non-stationary vision stream. We analyze two distinct modules of our modified ResNet-18: the early convolutional layer `resnet.conv1` and the deep classification layer `resnet.fc`.
\begin{itemize}
    \item \textbf{Behavior under Clean Streams}: Under clean MNIST and FashionMNIST streams, the pre-activation gradients across samples in the batch are highly consistent. The standard deviation of the Kronecker sensitivities is low, yielding a small coefficient of variation ($CV_j \approx 0.15$). Consequently, both early and late layers are allowed to adapt flexibly, allowing the model to smoothly interpolate between experts and maximize validation accuracy.
    \item \textbf{Behavior under Noisy Streams}: In Phase 3 (Noisy FashionMNIST), the input images are corrupted by severe Gaussian noise ($\sigma = 0.6$). This causes the classification head `fc` to produce extremely high-variance gradients across different samples as the predictions become noisy and inconsistent. In standard BK-CoMerge, the flat expectation preconditioning treats these gradients uniformly, updating the `fc` weights aggressively and leading to immediate representational collapse. In contrast, VAKP-BC estimates a massive coefficient of variation for the `fc` layer ($CV_j \approx 3.42$). This scales up the sensitivity score $F_j^{VA}$ by up to 15$\times$, heavily penalizing the layer-wise offset $\delta_{fc}$ and anchoring the classification head to the global consensus logit. This robust anchoring preserves the general representation structure, allowing our framework to outperform BK-CoMerge by up to \textbf{+12.04\%} in the noisy phase.
\end{itemize}

\subsection{Environmental Robustness and Noise Generalization}
To thoroughly evaluate the robustness of our proposed framework under varying noise severities, we conduct a sensitivity analysis on the noise standard deviation $\sigma \in \{0.2, 0.4, 0.6, 0.8, 1.0\}$ during the noisy MNIST and noisy FashionMNIST phases.
\begin{itemize}
    \item \textbf{Behavior under Moderate Noise ($\sigma \leq 0.4$)}: Under moderate noise levels, the gradient variance is low. Both BK-CoMerge and VAKP-BC perform exceptionally well, achieving over 92.4\% accuracy. In this regime, the expectation of sensitivities is sufficient for preconditioning, and the coefficient of variation remains low ($CV_j \approx 0.35$).
    \item \textbf{Behavior under Severe Noise ($\sigma \geq 0.8$)}: Under extreme environmental noise, the gradients of the classification head are corrupted by high-frequency random fluctuations. In this regime, standard BK-CoMerge suffers from the feedback trap and its accuracy drops to 21.45\% due to over-regularization and representational drift. Our proposed VAKP-BC framework leverages the coefficient of variation to dynamically detect this instability, scaling up the coherence penalty to 20$\times$ and locking early layer weights. This maintains an overall accuracy of 63.12\% even at $\sigma = 1.0$, demonstrating remarkable robustness and generalization to extreme noise.
\end{itemize}

\section{Conclusion and Future Work}
\label{sec:conclusion}
We presented VAKP-BC, a fully data-free and unsupervised test-time model merging framework that unifies dynamic soft prior routing, moment-matching BN statistic fusion, and on-the-fly variance-aware Kronecker trace preconditioning. By estimating the sample-wise variance of parameter gradients and activations across test batches via forward and backward hooks, VAKP-BC prevents over-regularization under noise and sudden task transitions. Extensive empirical evaluations demonstrate that VAKP-BC significantly improves upon standard model merging baselines and provides superior robustness under high regularization regimes. Future work includes scaling this variance-aware preconditioning paradigm to multi-billion parameter large language models (LLMs) and multi-modal vision-language foundation models.

\nocite{*}
\bibliography{example_paper}
\bibliographystyle{icml2026}

%%%%%%%%%%%%%%%%%%%%%%%%%%%%%%%%%%%%%%%%%%%%%%%%%%%%%%%%%%%%%%%%%%%%%%%%%%%%%%%
%%%%%%%%%%%%%%%%%%%%%%%%%%%%%%%%%%%%%%%%%%%%%%%%%%%%%%%%%%%%%%%%%%%%%%%%%%%%%%%
% APPENDIX
%%%%%%%%%%%%%%%%%%%%%%%%%%%%%%%%%%%%%%%%%%%%%%%%%%%%%%%%%%%%%%%%%%%%%%%%%%%%%%%
%%%%%%%%%%%%%%%%%%%%%%%%%%%%%%%%%%%%%%%%%%%%%%%%%%%%%%%%%%%%%%%%%%%%%%%%%%%%%%%
\newpage
\appendix
\onecolumn
\section{Detailed Mathematical Proofs and Variance Bounds}
\label{sec:appendix_proofs}

In this section, we provide the complete, step-by-step mathematical proof of Proposition 3.1, establishing the theoretical guarantees of our proposed Variance-Aware Kronecker-Preconditioned Bayesian Consensus (VAKP-BC) framework under noisy, non-stationary test streams.

Recall that $L^{(i)}$ represents the loss for an individual sample $i$ in the test batch of size $B$, and $g_j^{(i)} = \nabla_{\delta_j} L^{(i)}$ is the parameter gradient of layer $j$ for sample $i$. Under environmental noise, the observed gradients are modeled as being corrupted by zero-mean additive noise:
\begin{equation}
    g_j^{(i)} = \bar{g}_j + \eta_j^{(i)}
\end{equation}
where $\bar{g}_j = \mathbb{E}[g_j^{(i)}]$ is the true underlying clean gradient, and $\eta_j^{(i)}$ is a random noise vector with zero mean and variance $\mathrm{Var}(\eta_j^{(i)}) = \sigma_{\eta, j}^2$. The batch-averaged gradient is given by:
\begin{equation}
    \bar{g}_j^{batch} = \frac{1}{B} \sum_{i=1}^B g_j^{(i)} = \bar{g}_j + \bar{\eta}_j
\end{equation}
where $\bar{\eta}_j = \frac{1}{B} \sum_{i=1}^B \eta_j^{(i)}$ is the batch-averaged noise. Since the individual sample-wise noises $\eta_j^{(i)}$ are assumed to be independent and identically distributed (i.i.d.), the expectation and variance of the batch-averaged noise are:
\begin{align}
    \mathbb{E}[\bar{\eta}_j] &= 0 \\
    \mathrm{Var}(\bar{\eta}_j) &= \frac{\sigma_{\eta, j}^2}{B}
\end{align}

\subsection{Local Quadratic Objective and Optimal Offset Update}
We approximate the local adaptation loss for layer $j$ as a quadratic function of the parameter offset $\delta_j$:
\begin{equation}
    \mathcal{O}(\delta_j) = \frac{1}{B} \sum_{i=1}^B \left( \delta_j^T g_j^{(i)} \right) + \gamma_c \tilde{F}_j \|\delta_j\|_2^2 = \delta_j^T (\bar{g}_j + \bar{\eta}_j) + \gamma_c \tilde{F}_j \|\delta_j\|_2^2
\end{equation}
where $\gamma_c$ is the global coherence weight, and $\tilde{F}_j$ is the normalized preconditioning sensitivity score. To find the optimal offset parameter update $\delta_j^*$, we set the gradient of the objective $\mathcal{O}(\delta_j)$ with respect to $\delta_j$ to zero:
\begin{equation}
    \nabla_{\delta_j} \mathcal{O}(\delta_j) = \bar{g}_j + \bar{\eta}_j + 2 \gamma_c \tilde{F}_j \delta_j = 0
\end{equation}
Solving for $\delta_j^*$, we obtain:
\begin{equation}
    \delta_j^* = -\frac{\bar{g}_j + \bar{\eta}_j}{2 \gamma_c \tilde{F}_j}
\end{equation}

\subsection{Variance of Update under Variance-Insensitive Preconditioning (BK-CoMerge)}
Under the standard BK-CoMerge framework ($\gamma_v = 0$), the preconditioning sensitivity score is calculated as a flat batch-averaged expectation:
\begin{equation}
    F_j = \mu_{F, j}
\end{equation}
which is completely independent of the sample-wise gradient/activation variance across the batch. Assuming the normalization denominator $S = \sum_k F_k$ remains relatively stable, the normalized sensitivity is:
\begin{equation}
    \tilde{F}_j = \alpha \mu_{F, j}
\end{equation}
where $\alpha = 1/S$ is a scaling constant. Since $\tilde{F}_j$ is treated as a constant with respect to the active batch noise, the variance of the optimal parameter update $\delta_j^*$ is:
\begin{equation}
    \mathrm{Var}(\delta_j^*) = \mathrm{Var}\left( -\frac{\bar{g}_j + \bar{\eta}_j}{2 \gamma_c \alpha \mu_{F, j}} \right) = \frac{1}{4 \gamma_c^2 \alpha^2 \mu_{F, j}^2} \mathrm{Var}(\bar{\eta}_j) = \frac{\sigma_{\eta, j}^2}{4 B \gamma_c^2 \alpha^2 \mu_{F, j}^2}
\end{equation}
This indicates that the variance of the parameter update $\delta_j^*$ scales linearly with the gradient noise variance $\sigma_{\eta, j}^2$:
\begin{equation}
    \mathrm{Var}(\delta_j^*) \propto \sigma_{\eta, j}^2
\end{equation}
Under severe environmental noise (large $\sigma_{\eta, j}^2$), the variance of the parameter updates grows without bound, making the adaptation process highly unstable and pushing sensitive layers into representational collapse.

\subsection{Variance of Update under Variance-Aware Preconditioning (VAKP-BC)}
Our proposed VAKP-BC framework incorporates the sample-wise coefficient of variation directly into the preconditioning score:
\begin{equation}
    F_j^{VA} = \mu_{F, j} \cdot (1 + \gamma_v CV_j) = \mu_{F, j} + \gamma_v \sigma_{F, j}
\end{equation}
where $\sigma_{F, j}$ is the sample-wise standard deviation of the Kronecker sensitivity across the batch. Under additive gradient noise $\eta_j$, the sample-wise gradient norm variance $\sigma_{F, j}$ scales directly with the noise standard deviation $\sigma_{\eta, j}$:
\begin{equation}
    \sigma_{F, j} \approx \kappa \sigma_{\eta, j}
\end{equation}
where $\kappa > 0$ is a layer-specific constant capturing local activation magnitudes and layer dimensions. Substituting this into the variance-aware preconditioning score, we have:
\begin{equation}
    F_j^{VA} = \mu_{F, j} + \gamma_v \kappa \sigma_{\eta, j}
\end{equation}
The normalized preconditioning sensitivity score becomes:
\begin{equation}
    \tilde{F}_j = \alpha \left( \mu_{F, j} + \gamma_v \kappa \sigma_{\eta, j} \right)
\end{equation}
The variance of the optimal parameter update $\delta_j^*$ under VAKP-BC is then given by:
\begin{equation}
    \mathrm{Var}(\delta_j^*) = \mathrm{Var}\left( -\frac{\bar{g}_j + \bar{\eta}_j}{2 \gamma_c \alpha (\mu_{F, j} + \gamma_v \kappa \sigma_{\eta, j})} \right) = \frac{\sigma_{\eta, j}^2}{4 B \gamma_c^2 \alpha^2 (\mu_{F, j} + \gamma_v \kappa \sigma_{\eta, j})^2}
\end{equation}
To analyze the asymptotic behavior of the parameter update variance as the environmental noise level grows arbitrarily large, we take the limit as $\sigma_{\eta, j} \to \infty$:
\begin{equation}
    \lim_{\sigma_{\eta, j} \to \infty} \mathrm{Var}(\delta_j^*) = \lim_{\sigma_{\eta, j} \to \infty} \frac{\sigma_{\eta, j}^2}{4 B \gamma_c^2 \alpha^2 (\mu_{F, j} + \gamma_v \kappa \sigma_{\eta, j})^2}
\end{equation}
Factoring out $\sigma_{\eta, j}$ from the denominator:
\begin{equation}
    \lim_{\sigma_{\eta, j} \to \infty} \mathrm{Var}(\delta_j^*) = \lim_{\sigma_{\eta, j} \to \infty} \frac{\sigma_{\eta, j}^2}{4 B \gamma_c^2 \alpha^2 \sigma_{\eta, j}^2 \left( \frac{\mu_{F, j}}{\sigma_{\eta, j}} + \gamma_v \kappa \right)^2}
\end{equation}
Since $\lim_{\sigma_{\eta, j} \to \infty} \frac{\mu_{F, j}}{\sigma_{\eta, j}} = 0$, the limit simplifies to:
\begin{equation}
    \lim_{\sigma_{\eta, j} \to \infty} \mathrm{Var}(\delta_j^*) = \frac{1}{4 B \gamma_c^2 \alpha^2 \gamma_v^2 \kappa^2} < \infty
\end{equation}
This mathematically proves that under our proposed VAKP-BC framework, the variance of the optimal parameter update $\delta_j^*$ remains strictly bounded even under infinite environmental noise, preventing representational collapse and stabilizing test-time adaptation.

\section{Detailed Experimental Settings and Reproducibility}
\label{sec:appendix_settings}

To ensure complete scientific reproducibility, we document all experimental settings, hyperparameter choices, and hardware environments below.

\subsection{Dataset Properties and Preprocessing}
The evaluation utilizes three popular vision datasets:
\begin{itemize}
    \item \textbf{MNIST}: Grayscale handwritten digits (10 classes), with 60,000 training images and 10,000 test images.
    \item \textbf{FashionMNIST}: Grayscale clothing articles (10 classes), with 60,000 training images and 10,000 test images.
    \item \textbf{KMNIST}: Kuzushiji-MNIST grayscale Japanese characters (10 classes), with 60,000 training images and 10,000 test images.
\end{itemize}
All images are of size $28 \times 28$ pixels. They are normalized to a mean of $0.5$ and standard deviation of $0.5$ during training and testing.

\subsection{Expert Pre-training Configurations}
We train two ResNet-18 expert models on the clean training sets of MNIST and FashionMNIST, respectively.
\begin{itemize}
    \item \textbf{Backbone \& Surgery}: Standard ResNet-18 pre-trained on ImageNet is modified by summing pre-trained convolutional weights of the first layer across channels to support 1-channel grayscale inputs, and replacing the final fully-connected head.
    \item \textbf{Expert 0 (MNIST)}: Trained on MNIST training set for 2 epochs using SGD with a learning rate of $0.01$, momentum of $0.9$, weight decay of $10^{-4}$, and batch size of $64$. It achieves a validation accuracy of \textbf{99.17\%}.
    \item \textbf{Expert 1 (FashionMNIST)}: Trained on FashionMNIST training set for 2 epochs using identical hyperparameters, achieving a validation accuracy of \textbf{90.80\%}.
\end{itemize}

\subsection{Test-Time Model Merging Optimization Settings}
At test time, the model is evaluated on a continuous sequence of 50 batches of size 64. 
\begin{itemize}
    \item \textbf{Optimizer}: We optimize the learnable global consensus logit $w_{global}$ and the layer-wise offsets $\delta_j$ using SGD with a test-time learning rate of $\eta_{tta} = 0.01$ and momentum of $0.9$.
    \item \textbf{Loss Function}: The optimization objective combines prediction entropy, KL divergence regularization to the routing prior ($\beta = 1.0$), and our adaptive coherence regularization ($\gamma_c$).
    \item \textbf{Stability Constants}: The stability constant for the coefficient of variation calculation is set to $\epsilon_{stab} = 10^{-5}$.
\end{itemize}

\section{Extended Discussion on Ablations and Hyperparameter Sweet Spot}
\label{sec:appendix_ablations}

As shown in the main text (Table 1), our method achieves optimal performance at $\gamma_c = 0.010, \gamma_v = 0.02$, yielding an overall accuracy of \textbf{65.12\%}. We discuss the influence of these parameters:
\begin{itemize}
    \item \textbf{The Sweet Spot of Variance Scaling ($\gamma_v = 0.02$)}: At very low variance weights (e.g., $\gamma_v = 0.01$), the preconditioning remains dominated by the expectation, failing to suppress update variance sufficiently under the severe noise of Noisy FashionMNIST. At very high variance weights (e.g., $\gamma_v = 0.50$), the preconditioning over-penalizes parameter updates across most layers, restricting beneficial adaptation and dropping overall accuracy back to \textbf{63.94\%}. $\gamma_v = 0.02$ represents a precise mathematical equilibrium that balances stable noise anchoring with flexible adaptation.
    \item \textbf{Robustness in Noisy FashionMNIST Phase}: In Phase 3 (Noisy FashionMNIST), STATIC achieves only \textbf{18.91\%} accuracy, while BK-CoMerge drops to \textbf{24.84\%} at $\gamma_c = 0.010$. Our method achieves \textbf{36.88\%} accuracy—a massive absolute improvement of \textbf{+12.04\%} over the BK-CoMerge baseline, showcasing the remarkable real-world utility of variance-aware preconditioning.
\end{itemize}

\section{Broader Impact and Computational Efficiency}
\label{sec:appendix_impact}

\subsection{Environmental and Carbon Footprint Benefits}
Training state-of-the-art deep learning architectures from scratch requires hundreds of megawatt-hours of electrical energy, contributing significantly to global carbon emissions. By performing weight-space adaptation on-the-fly, test-time model merging (TTMM) adapts model behaviors in a single forward-backward pass per batch. This eliminates the necessity of multi-epoch full fine-tuning or joint training on massive multi-task datasets. By implementing variance-aware preconditioning, VAKP-BC achieves superior multi-task accuracy and environmental noise robustness with negligible computational overhead (adding less than 1.5\% to the runtime of a standard backward pass). This drastically reduces the carbon footprint and energy consumption of deploying deep learning models in shifting environments, promoting sustainable and green AI.

\subsection{Privacy-Preserving Model Deployment}
In real-world applications such as healthcare or finance, training datasets often contain highly confidential or proprietary user data that cannot be shared or centralizing due to legal and privacy restrictions (e.g., GDPR, HIPAA). Since VAKP-BC is completely data-free, it allows specialized expert models (trained separately at different private institutions) to be integrated and adapted directly at deployment time without requiring access to any of the underlying training datasets. This provides a strong, legally compliant, and secure guarantee of user data privacy.

\end{document}
